\documentclass[conference]{IEEEtran}
\usepackage{graphicx,booktabs,hyperref}
\graphicspath{{figures/}}

\begin{document}
\title{Measuring Wi-Fi and Cellular Performance Along a University Route}
\author{\IEEEauthorblockN{K.K.S. Kapuruge (220301V) \quad W.S.S. Jayawardhana (220282K)}
\IEEEauthorblockA{CS4473 Mobile Computing --- Group 11, Option F: Mobile Network Measurement}}
\maketitle

\begin{abstract}
% 3-4 sentences: what we measured, how, the main finding (with one number), one limitation.
\end{abstract}

\section{Introduction}
% Problem: network quality changes with place and time; mobile users move between Wi-Fi and cellular.
% Why it matters for mobile computing: intermittent connectivity, variable signal quality.
% Research question (from the proposal) and expectation.
% At least two references, e.g. Android WifiInfo / TelephonyManager docs \cite{wifiinfo,telephony}.

\section{Prototype and Method}
\subsection{Prototype}
% Block diagram (Fig.~\ref{fig:block}), APIs: WifiManager/WifiInfo via NetworkCallback, TelephonyManager,
% LocationManager, ConnectivityManager.NetworkCallback, HttpURLConnection, foreground service (location).
% Probe: 3 x /ping (median) + 256 000-byte download every 10 s and on waypoint taps; skip rule.
% Server: laptop + fixed ngrok URL.
\subsection{Experiment}
% Route GTN Lab -> Library, P1-P5, 40 s pause. Factors: network (Wi-Fi/cellular) x time slot
% (08:00-09:00 / 12:30-13:30). 3 walks per slot, both phones; role swap. Phones: 2 x Pixel 7 Pro,
% Android 16, Mobitel. Held constant: ... Ground truth: waypoint taps, RSSI vs Settings, bytes vs server log.

\section{Results}
% Table~\ref{tab:points}: mean +- SD per point (from analysis/output/point_stats.md).
% Fig.~\ref{fig:route}: latency / speed along the route.
% Fig.~\ref{fig:scatter}: signal vs latency and speed, with r.
% Drops per walk and where (drops.md). Peak vs off-peak. Did results match the expectation?

\begin{figure}[t]
\centering
\includegraphics[width=\columnwidth]{fig_route_dl_mbps.png}
\caption{Download speed at P1--P5 (mean $\pm$ SD over 3 walks).}
\label{fig:route}
\end{figure}

\begin{figure}[t]
\centering
\includegraphics[width=\columnwidth]{fig_scatter.png}
\caption{Signal strength vs latency and download speed.}
\label{fig:scatter}
\end{figure}

\section{Discussion}
% Mobile computing challenges met and how we handled them: Android 16 Wi-Fi info redaction,
% foreground-service rules, Doze, permissions, ngrok warning page, clock sync between phones.
% Limitations: one route, one operator, n = 3 walks, ngrok detour adds latency, 250 KB file
% underestimates throughput (TCP slow start), campus Wi-Fi load not controlled.

\section{Contributions}
% K.K.S. Kapuruge: ...
% W.S.S. Jayawardhana: ...
% AI use: see AI_USAGE_LOG.md in the data package.

\begin{thebibliography}{9}
\bibitem{wifiinfo} Android Developers, ``WifiInfo,'' \url{https://developer.android.com/reference/android/net/wifi/WifiInfo}.
\bibitem{telephony} Android Developers, ``TelephonyManager,'' \url{https://developer.android.com/reference/android/telephony/TelephonyManager}.
\end{thebibliography}
\end{document}
